\selectlanguage{english}
\chapter[Extension host]{VS Code extension host and auxiliary modules}
\label{chap4:extension}
\begin{chapquote}{Microsoft Documentation}
``Extensions can interact with the editor through the VS Code Extension API.''
\end{chapquote}

\section{Purpose and responsibilities}
\label{chap4:purpose}

\texttt{extension.ts} is the \textbf{activation entry point} of the DI Plugin. The \texttt{activate(context)} function registers all subscriptions; \texttt{deactivate()} disposes diagnostics. Responsibilities span the full user journey: from opening a C\# file and seeing constructor hints, through workspace-wide missing-registration reports, to writing generated factories and merging \texttt{services.Add*} lines into a composition root.

The host module is intentionally large because it encodes domain heuristics specific to \texttt{Microsoft.Extensions.DependencyInjection} and ASP.NET Core idioms. Two auxiliary modules isolate edit construction to keep the main file maintainable:
\begin{itemize}
	\item \texttt{registrationInsert.ts} --- discovery of insertion anchors, deduplication, parsing generated \texttt{AppBuilder.cs}, lifetime normalization;
	\item \texttt{extractInterface.ts} --- multi-file \texttt{WorkspaceEdit} for concrete-parameter quick fixes.
\end{itemize}

This chapter documents host data structures, the workspace aggregation algorithm, each contributed command workflow, the VS Code Extension API mapping (Section~\ref{chap4:vscode-api}), diagnostic and code-action behavior, design heuristics, and limitations. Empirical outcomes appear in Chapter~\ref{chap5:evaluation}; screenshot guidance appears in Chapter~\ref{chap6:illustrations}.

\section{Key data structures and interfaces}
\label{chap4:data-structures}

Beyond \texttt{WorkspaceDiSummary} (Chapter~\ref{chap2:data-structures}), the host defines supporting types summarized in Table~\ref{tab:host-types}.

\begin{table}[htbp]
\centering
\caption{Host-specific structures}
\label{tab:host-types}
\begin{tabular}{|l|p{8.5cm}|}
\hline
\textbf{Symbol} & \textbf{Role} \\
\hline
\texttt{WorkspaceCSharpFile} & \texttt{uri}, \texttt{source}, optional \texttt{document} reference. \\
\hline
\texttt{GeneratedFile} & Target path under \texttt{GeneratedDI/} plus file body. \\
\hline
\texttt{RegistrationInsertTarget} & \texttt{uri}, \texttt{insertPosition}, \texttt{indent}, \texttt{label}. \\
\hline
\texttt{MARKER\_SERVICE\_INTERFACES} & Set excluding \texttt{IDisposable}, \texttt{ICloneable}, etc. from interface$\rightarrow$impl suggestions. \\
\hline
\texttt{DiCodeActionProvider} & \texttt{CodeActionProvider} for DI-related diagnostics. \\
\hline
\end{tabular}
\end{table}

\subsection{Implicit framework dependencies}

Function \texttt{isImplicitFrameworkDependency(typeName)} returns true for types the .NET host registers without explicit application code, including \texttt{IServiceProvider}, \texttt{IServiceScopeFactory}, \texttt{IHostEnvironment}, \texttt{ILogger}, \texttt{IConfiguration}, and generic patterns \texttt{ILogger<}, \texttt{IOptions<}, \texttt{IOptionsSnapshot<}. This prevents workspace reports from listing hundreds of false missing types when factories inject \texttt{IServiceProvider}.

\subsection{Non-injectable class detection}

Classes whose base type matches \texttt{Profile}, \texttt{Migration}, \texttt{Controller}, \texttt{PageModel}, \texttt{DbContext} (when used as inheritance marker), etc., are added to \texttt{nonInjectableClassNames}. AutoMapper profiles and EF migrations should be discovered by assembly scanning, not registered as \texttt{AddScoped<Profile>}. MealPlanner contains many \texttt{Profile} subclasses; excluding them was required to avoid absurd registration suggestions.

\section{Workspace aggregation algorithm}
\label{chap4:workspace-algo}

\texttt{buildWorkspaceSummary()} is the computational core for all workspace commands. Pseudocode:

\begin{algorithmic}[1]
\State $files \gets$ find up to 200 \texttt{*.cs} excluding build folders
\For{each $f \in files$}
\State $analysis \gets \texttt{analyzeCSharp}(f.source)$ with try/catch
\State merge $analysis.constructors$ with path metadata
\State $registeredTypes.\texttt{addAll}(\texttt{collectExtendedRegistrationTypes}(f.source))$
\State update $interfaceToImpl$, $typeToNamespace$, flags from regex
\EndFor
\State return \texttt{WorkspaceDiSummary}
\end{algorithmic}

\subsection{\texttt{collectExtendedRegistrationTypes}}

This function extends plain \texttt{AddSingleton|Scoped|Transient<T>} scanning:
\begin{itemize}
	\item \texttt{AddDbContext<T>} registers \texttt{T} and \texttt{DbContextOptions<T>};
	\item \texttt{AddAutoMapper(...)} registers \texttt{IMapper};
	\item \texttt{AddIdentity<TUser,TRole>} and \texttt{AddIdentityCore<TUser>} register Identity managers and stores;
	\item \texttt{AddRoles<T>} / role managers add \texttt{RoleManager<T>};
	\item \texttt{typeof(IContract)}, \texttt{typeof(Impl)} pairs inside \texttt{new ServiceDescriptor(...)} (EventPlanner style).
\end{itemize}

MealPlanner validation (Section~\ref{chap5:mealplanner}) depended on this extension: without it, \texttt{IMapper} and Identity types appeared missing despite \texttt{AddAutoMapper} and \texttt{AddIdentity} calls.

\subsection{Interface implementation map}

\texttt{collectInterfaceImplementations} applies regex akin to \texttt{public class Impl : IFoo, IBar} and records \texttt{interfaceToImpl.set("IFoo", ["Impl"])}. Multiple implementations of the same interface append to the array; registration suggestion picks the first or prompts implicitly through ordering. Generic interfaces (\texttt{IBaseService<,>}) are partially supported via patterns in source.

\subsection{Missing type computation (workspace)}

For each constructor parameter type \(p\) in the aggregated list:
\begin{itemize}
	\item skip primitives and implicit framework types;
	\item skip if \(p \in registeredTypes\);
	\item skip if \(p\) is a non-injectable class name;
	\item else report as missing for workspace analyze output and suggestion builders.
\end{itemize}

\section{Extension commands and workflows}
\label{chap4:commands}

Table~\ref{tab:commands} lists palette commands from \texttt{package.json}. Two additional command IDs exist for code actions only.

\begin{table}[htbp]
\centering
\caption{Contributed commands and workflows}
\label{tab:commands}
\small
\begin{tabular}{|p{3.8cm}|p{9cm}|}
\hline
\textbf{Command ID} & \textbf{Workflow} \\
\hline
\texttt{di-plugin.analyzeFile} & Active editor C\# $\rightarrow$ \texttt{analyzeCSharp} $\rightarrow$ Output report; refresh diagnostics. \\
\hline
\texttt{di-plugin.suggestRegistrations} & Current file registration lines via \texttt{buildRegistrationSuggestions}. \\
\hline
\texttt{di-plugin.analyzeWorkspace} & \texttt{buildWorkspaceSummary} $\rightarrow$ missing types, constructors grouped by file. \\
\hline
\texttt{di-plugin.suggestWorkspaceRegistrations} & \texttt{buildWorkspaceRegistrationSuggestions} with \texttt{preferredLifetimeForWorkspace}. \\
\hline
\texttt{di-plugin.suggestFactories} & File factories; if empty, fallback workspace factory snippets. \\
\hline
\texttt{di-plugin.suggestWorkspaceFactories} & \texttt{buildWorkspaceFactorySuggestions} --- factory \texttt{.cs} text only. \\
\hline
\texttt{di-plugin.generateWorkspaceFactories} & Write \texttt{GeneratedDI/*Factory.cs} + \texttt{AppBuilder.cs}. \\
\hline
\texttt{di-plugin.insertWorkspaceRegistrations} & Suggest then \texttt{insertRegistrationsIntoWorkspace}. \\
\hline
\texttt{di-plugin.mergeGeneratedIntoCompositionRoot} & Parse \texttt{AppBuilder} registrations $\rightarrow$ insert into root. \\
\hline
\end{tabular}
\end{table}

\subsection{DI: Analyze current file}

\textbf{User goal:} understand constructor dependencies in the file being edited.

\textbf{Steps:} resolve active text editor; verify language C\#; call \texttt{analyzeCSharp}; format constructors with parameter lists; list concrete-type, cycle, and file-local missing issues; show errors including parse limit message.

\textbf{Sample result:} On \texttt{CheckoutSample/Services.cs}, output lists \texttt{CheckoutCoordinator} with four interface parameters and \texttt{PaymentCaptureService} with two---establishing factory eligibility without workspace scan.

\subsection{DI: Analyze workspace}

\textbf{User goal:} solution-wide checklist of missing registrations and constructor inventory.

\textbf{Steps:} \texttt{buildWorkspaceSummary}; diff constructor parameter types against \texttt{registeredTypes}; print grouped report to Output channel.

\textbf{Sample result:} On \texttt{BuilderFactorySample}, reports missing \texttt{IEmailSender} and \texttt{OrderNotificationService} (or factory) relative to partial \texttt{CompositionRoot}.

\subsection{DI: Suggest registrations (file and workspace)}

\texttt{buildWorkspaceRegistrationSuggestions} emits lines:
\begin{itemize}
	\item \texttt{services.Add\{Lifetime\}<IInterface, Impl>()} for each unregistered interface key in \texttt{interfaceToImpl};
	\item \texttt{services.Add\{Lifetime\}<Concrete>()} for concrete classes with constructors that are not factory targets, not registered, not excluded by MVC/test/non-injectable rules.
\end{itemize}

\texttt{preferredLifetimeForWorkspace} algorithm:
\begin{enumerate}
	\item If composition root file exists, count \texttt{AddSingleton}, \texttt{AddScoped}, \texttt{AddTransient} occurrences; pick majority.
	\item Else if \texttt{aspNetMvcHost}, default \texttt{Scoped}.
	\item Else default \texttt{Transient} (console apps like BaseSample).
\end{enumerate}

\textbf{Regression fix:} \texttt{buildRegistrationLineForMissingType} maps concrete \texttt{SmtpEmailSender} to \texttt{ISmtpEmailSender} pair when suggesting, avoiding duplicate ``already registered'' false negatives when interface form is required.

\subsection{DI: Suggest factories (file and workspace)}

\texttt{shouldGenerateFactory(ctor, interfaceToImpl)}:
\[
|\text{params}| \geq 3 \quad \land \quad \forall p.\; \texttt{interfaceToImpl.has}(p.\text{type})
\]

\textbf{CheckoutSample:} \texttt{CheckoutCoordinator} qualifies; \texttt{PaymentCaptureService} fails on count.

\textbf{MultiFactorySample:} \texttt{ImportPipeline} and \texttt{ExportPipeline} both qualify; \texttt{JobScheduler} fails.

\texttt{buildWorkspaceFactorySuggestions} returns only sealed factory class templates with \texttt{Create()} calling \texttt{GetRequiredService<T>()} for each dependency---not the full AppBuilder.

\subsection{DI: Generate builders/factories files (workspace)}

Writes to \texttt{GeneratedDI/}:
\begin{itemize}
	\item one \texttt{*Factory.cs} per eligible class;
	\item \texttt{AppBuilder.cs} static class with \texttt{Build()} method.
\end{itemize}

\texttt{buildAppBuilderFileContent} rules (refined):
\begin{itemize}
	\item register \textbf{factory types} for factory-backed services, not the concrete orchestrator type;
	\item include interface$\rightarrow$implementation lines required by factory \texttt{Create()} even if already present in manual \texttt{CompositionRoot} (standalone checklist semantics);
	\item skip marker interfaces and non-injectable classes;
	\item add \texttt{using} statements from \texttt{typeToNamespace}.
\end{itemize}

\textbf{MultiFactorySample committed output} registers \texttt{ImportPipelineFactory} and \texttt{ExportPipelineFactory} plus dependency interfaces---documented in Chapter~\ref{chap5:multifactory}.

\subsection{DI: Insert workspace registrations}

Pipeline: \texttt{buildWorkspaceRegistrationSuggestions} $\rightarrow$ \texttt{findRegistrationInsertTarget} $\rightarrow$ \texttt{filterNewRegistrationLines} $\rightarrow$ \texttt{WorkspaceEdit.insert} at computed position $\rightarrow$ \texttt{applyEdit}.

User sees success notification with target label (\texttt{CompositionRoot}, etc.).

\subsection{DI: Merge GeneratedDI into CompositionRoot}

Reads \texttt{GeneratedDI/AppBuilder.cs}, extracts \texttt{services.Add*} lines via line-by-line parser (\texttt{parseRegistrationsFromAppBuilder}), deduplicates against existing root content, inserts using same anchor as insert command.

\textbf{Regression:} early regex failed on \texttt{AddScoped<T>()} without interface pair; line-based parser fixed merge for BuilderFactorySample.

\subsection{Internal commands (code actions)}

\texttt{di-plugin.insertRegistrationForType} --- triggered from lightbulb on missing-registration diagnostic; passes type name argument.

\texttt{di-plugin.extractInterfaceAndRegister} --- runs \texttt{extractInterface.ts} workflow then inserts registration.

\section{Registration insert module (\texttt{registrationInsert.ts})}
\label{chap4:registration-insert}

\subsection{Purpose}

Encapsulates brittle text-editing rules so \texttt{extension.ts} does not duplicate position logic across insert and merge.

\subsection{\texttt{findRegistrationInsertTarget}}

Priority order:
\begin{enumerate}
	\item file containing \texttt{class CompositionRoot};
	\item \texttt{public static IServiceCollection Add*Services(this IServiceCollection ...)};
	\item \texttt{void ConfigureServices(IServiceCollection services)} in startup style.
\end{enumerate}

Within the target file, search backward from \texttt{return services.BuildServiceProvider()} or \texttt{return services;}. Insert position is \textbf{start of that line's} zero-based offset, not the \texttt{return} token offset---preventing corrupted lines such as \texttt{re services.Add...turn}.

\subsection{\texttt{filterNewRegistrationLines}}

Normalizes whitespace and rejects lines whose substantive content already appears in the target file, reducing duplicate \texttt{AddScoped} on repeated merge.

\subsection{\texttt{parseRegistrationsFromAppBuilder}}

Splits \texttt{AppBuilder.cs} into lines; keeps lines matching \texttt{services.Add(Singleton|Scoped|Transient)}. Supports one-type-parameter and two-type-parameter generic forms.

\section{Extract interface module (\texttt{extractInterface.ts})}
\label{chap4:extract-interface}

\subsection{Workflow}

For concrete parameter type \texttt{SmtpEmailSender} in consumer class:
\begin{enumerate}
	\item derive interface name \texttt{ISmtpEmailSender} (prefix \texttt{I} + concrete name, or user convention);
	\item locate implementation file defining \texttt{SmtpEmailSender};
	\item create interface file under \texttt{Interfaces/} if that folder exists in workspace, else adjacent to implementation;
	\item add interface to class base list on implementation;
	\item replace parameter type in consumer constructor with interface;
	\item call \texttt{insertRegistrationsIntoWorkspace} with \texttt{AddScoped<ISmtpEmailSender, SmtpEmailSender>()}.
\end{enumerate}

\subsection{Limitations}

Generated interfaces are \textbf{empty}---members are not copied from implementation. Developer must move method signatures manually. This is acceptable for thesis prototype scope but noted in evaluation threats.

\section{VS Code Extension API usage in detail}
\label{chap4:vscode-api}

This section documents how the host maps architectural responsibilities onto concrete Extension API types~\cite{vscode-api}. Understanding these mappings explains activation behavior and suggests where screenshots should be taken for the thesis (Chapter~\ref{chap6:illustrations}).

\subsection{Activation and subscriptions}

\texttt{activate(context: vscode.ExtensionContext)} registers all listeners. Each subscription is pushed to \texttt{context.subscriptions} so VS Code disposes them when the extension unloads. \texttt{deactivate()} clears the diagnostic collection explicitly. Activation event \texttt{onLanguage:csharp} delays loading until a C\# file exists, reducing startup cost in non-.NET workspaces.

\subsection{Document and workspace events}

\begin{itemize}
	\item \texttt{workspace.onDidOpenTextDocument} / \texttt{onDidChangeTextDocument} --- if \texttt{document.languageId === 'csharp'}, call \texttt{analyzeCSharpDocument} to refresh diagnostics.
	\item \texttt{workspace.onDidCloseTextDocument} --- remove diagnostics for closed URI to avoid stale entries.
	\item \texttt{workspace.findFiles(include, exclude, maxResults)} --- workspace scan with glob \texttt{**/*.cs} and exclusion of \texttt{**/bin/**}, \texttt{**/obj/**}; \texttt{maxResults = 200}.
	\item \texttt{workspace.fs.readFile(uri)} --- returns \texttt{Uint8Array} decoded to string for analysis; failures skip file silently.
\end{itemize}

\subsection{Diagnostics API}

\texttt{languages.createDiagnosticCollection('di-plugin')} returns a mutable collection keyed by document URI. Each \texttt{Diagnostic} includes:
\begin{itemize}
	\item \texttt{range}: \texttt{new vscode.Range(startLine, startChar, endLine, endChar)} from constructor/param indices;
	\item \texttt{message}: human-readable dependency list or issue description;
	\item \texttt{severity}: \texttt{DiagnosticSeverity.Information} for constructor info, \texttt{Warning} for smells;
	\item \texttt{source}: \texttt{'di-plugin'} groups entries in Problems panel.
\end{itemize}

Parse-limit errors use range \texttt{(0,0)} at document start so the user sees why AST features are disabled for EF snapshots.

\subsection{Commands API}

\texttt{commands.registerCommand('di-plugin.analyzeWorkspace', async () => ...)} binds palette titles from \texttt{package.json} \texttt{contributes.commands}. Commands are async functions: they await file reads and \texttt{applyEdit}. Internal commands (\texttt{insertRegistrationForType}) accept \texttt{typeName} via \texttt{args} from code actions.

\subsection{Code actions API}

\texttt{languages.registerCodeActionsProvider} with \texttt{DiCodeActionProvider} implements \texttt{provideCodeActions(document, range, context)}. It inspects \texttt{context.diagnostics} filtered by \texttt{source === 'di-plugin'} and returns \texttt{CodeAction} objects with \texttt{kind: QuickFix} and \texttt{command} property (title + command id + arguments). Using commands rather than inline \texttt{WorkspaceEdit} in the provider ensures a single undo step and consistent error handling in the command handler.

\subsection{Workspace edit and file system API}

\texttt{insertRegistrationsIntoWorkspace} builds \texttt{WorkspaceEdit}:
\begin{itemize}
	\item \texttt{edit.insert(uri, position, text)} --- inserts registration lines at composition-root anchor;
	\item \texttt{workspace.applyEdit(edit)} --- applies atomically; returns boolean success.
\end{itemize}

\texttt{generateWorkspaceFactories} uses \texttt{workspace.fs.writeFile(uri, Buffer.from(content))} to create \texttt{GeneratedDI/} files. Directory creation uses \texttt{workspace.fs.createDirectory} when missing.

\subsection{Window and output API}

\texttt{window.createOutputChannel('DI Plugin')} receives append-only reports from analyze/suggest commands---useful when Problems panel diagnostics are insufficient (full workspace listing). \texttt{window.showInformationMessage} confirms successful insert/merge with composition-root label.

\subsection{Position conversion helper (conceptual)}

The host maps tree-sitter \texttt{startIndex} to \texttt{document.positionAt(offset)} when available, which accounts for UTF-16 in VS Code documents. This bridge is the critical link between the analyzer chapter and editor UX; misalignment would show squiggles on wrong tokens.

\section{Diagnostics and quick fixes}
\label{chap4:diagnostics}

\texttt{analyzeCSharpDocument(document)} clears and repopulates the diagnostic collection. Source attribute \texttt{di-plugin} groups entries in Problems panel.

\texttt{DiCodeActionProvider.provideCodeActions} returns actions when diagnostics match:
\begin{itemize}
	\item \textbf{Insert registration} --- local \texttt{ConfigureServices} when file matches startup pattern; else command to workspace insert with type argument;
	\item \textbf{Register interface} --- suggest \texttt{Add*<I, Concrete>} when \texttt{interfaceToImpl} has mapping;
	\item \textbf{Extract interface and register} --- full \texttt{extractInterfaceAndRegister} for concrete-type issues.
\end{itemize}

Code actions use \texttt{Command} with \texttt{arguments} rather than immediate \texttt{WorkspaceEdit} in provider, so execution happens after user selection with proper undo stack.

\section{Heuristics and design decisions}
\label{chap4:heuristics}

\begin{itemize}
	\item \textbf{Regex + AST hybrid} --- AST for constructors; regex for registrations and inheritance (pragmatic for ASP.NET extension methods scattered in large files).
	\item \textbf{Marker interfaces excluded} --- \texttt{IDisposable} on \texttt{UnitOfWork} must not produce \texttt{AddScoped<IDisposable, UnitOfWork>}.
	\item \textbf{Factory vs.\ registration UX split} --- reduces duplicate output; documented in README command table.
	\item \textbf{200-file cap} --- protects against freezing VS Code on accidental open of huge monorepos; trade-off documented.
	\item \textbf{Interface naming convention} --- \texttt{I + PascalCase} required for \texttt{collectInterfaceImplementations}; deviations are missed.
\end{itemize}

\section{Known limitations}
\label{chap4:limitations}

\begin{itemize}
	\item Workspace cap may omit files in solutions with more than 200 application \texttt{.cs} files.
	\item Lifetime inference is majority-vote, not architecture-aware (singleton cache in scoped service mistakes possible).
	\item Extract interface produces empty interfaces.
	\item \texttt{insertRegistrationForType} and \texttt{extractInterfaceAndRegister} are not in \texttt{package.json} contributes menu---only code actions.
	\item Generated \texttt{AppBuilder} may duplicate manual root until user regenerates after rule changes.
	\item No grouped undo title across multi-file extract beyond default VS Code behavior.
\end{itemize}

\subsection{Composition root excerpt (BuilderFactorySample)}

The partial wiring below (from \texttt{CompositionRoot.cs}) is intentional test input: the plugin must propose missing \texttt{IEmailSender} and notification/factory registration without duplicating existing lines.

\begin{lstlisting}[caption={Partial CompositionRoot (test input)}]
services.AddScoped<ILogger, ConsoleLogger>();
services.AddScoped<IClock, SystemClock>();
services.AddScoped<ITemplateEngine, BasicTemplateEngine>();
// Missing: IEmailSender, OrderNotificationService / factory
return services.BuildServiceProvider();
\end{lstlisting}

After \textbf{DI: Insert workspace registrations} and optional merge from \texttt{GeneratedDI/AppBuilder.cs}, the root should contain additional \texttt{AddScoped} lines immediately above the \texttt{return}, with indentation matching the surrounding method body.

\subsection{Command interaction diagram}

Workspace commands are not independent pipelines: they share \texttt{buildWorkspaceSummary()}. Figure~\ref{fig:command-deps} shows dependencies. Any improvement to \texttt{collectExtendedRegistrationTypes} benefits analyze, suggest, insert, merge, and generate alike.

\begin{figure}[htbp]
\centering
\fbox{\parbox{0.9\textwidth}{
\small
\texttt{buildWorkspaceSummary} $\rightarrow$ \texttt{analyzeWorkspace}\\
$\rightarrow$ \texttt{suggestWorkspaceRegistrations} $\rightarrow$ \texttt{insertWorkspaceRegistrations}\\
$\rightarrow$ \texttt{suggestWorkspaceFactories} $\rightarrow$ \texttt{generateWorkspaceFactories}\\
$\rightarrow$ \texttt{mergeGeneratedIntoCompositionRoot}
}}
\caption{Shared workspace summary across commands}
\label{fig:command-deps}
\end{figure}

The host module ties analyzer capabilities to user-visible outcomes validated sample-by-sample in Chapter~\ref{chap5:evaluation}.
